% frontmatter/guide.tex — 导读
\chapter*{导读：如何阅读本书}
\markboth{导读}{导读}

\section*{全书结构}

本书共 10 章，按照从底层到上层的顺序组织：

\begin{itemize}
  \item 第 1 章讲构建系统，建立工程基线认知
  \item 第 2--3 章讲核心执行路径：入口、查询、工具
  \item 第 4 章讲终端界面的 React 化实现
  \item 第 5 章讲上下文工程——AI 应用中最关键的工程问题之一
  \item 第 6 章讲多任务与多代理运行时
  \item 第 7 章讲扩展系统的多层架构
  \item 第 8 章讲安全与权限模型
  \item 第 9 章讲性能与可靠性
  \item 第 10 章做跨章节的设计哲学总结
\end{itemize}

\section*{阅读路径建议}

顺序阅读是最完整的路径，每章都为后续章节建立概念基础。

如果你对特定主题更感兴趣：

\begin{itemize}
  \item 想理解 AI 工具调用机制：第 2 章 $\rightarrow$ 第 3 章
  \item 想理解上下文管理与压缩：第 5 章（需要第 2 章基础）
  \item 想理解安全模型：第 8 章（需要第 3 章基础）
  \item 想理解扩展架构：第 7 章（需要第 3 章基础）
  \item 想理解多代理协作：第 6 章（需要第 2、3 章基础）
\end{itemize}

\section*{术语约定}

本书的核心术语在术语表（附录 A）中统一定义。正文中术语首次出现时会标注英文原名，后续使用统一叫法。

源码中的标识符（函数名、类型名、文件路径）保留英文原名，使用等宽字体排版。

\section*{方法论边界}

本书的分析方法以静态源码阅读为主。所有技术结论都尽可能追溯到具体的文件、函数和类型定义。

对于静态阅读难以确定的机制（如运行时状态转换、异步调度顺序），我们会通过动态验证补充——但会明确标注哪些结论来自运行观测，哪些来自纯静态分析。

基于源码结构的合理推断会与确定性事实区分。我们不会把"可能是这样"写成"就是这样"。

\section*{版本基线}

本书基于 Claude Code 2.1.88 版本。Claude Code 是一个快速迭代的产品，后续版本的实现细节可能发生变化。本书描述的是这个版本的工程快照，而非永恒不变的设计。
